\section{Observations: Prompt Precision, Human Oversight, Architectural Limitations}

Across the features covered so far, three observations recur often enough to be treated as general findings about Part 1 rather than as properties of any single feature.

The first is that the quality of the generated output tracked the precision of the prompt closely and directly. The frontend illustrates this at the level of visual detail: the closer a prompt specified what was wanted --- layout, spacing, specific colors and fonts --- the closer the generated result matched the intended design, while looser prompts left more for the agent to guess and produced results that needed further rounds of correction. The backend illustrates the same point at the level of missing requirements rather than visual detail: the missing email field, the missing uniqueness constraint, and the generic validation messages were not failures of the agent's reasoning so much as direct consequences of a prompt that had never specified any of the three. In both cases, the gap between what was generated and what was actually wanted traces back to what the prompt did or did not say, not to a limitation in what the agent was capable of producing when asked precisely.

The second is that human oversight was present but shallow for most of Part 1, and this is what allowed the architectural problems discussed earlier to accumulate rather than being caught early. A standing rule prevented the agent from modifying any existing file without asking first, but in practice this meant changes were generally accepted once they worked, without being checked against the system they were being added to --- the service layer, the growing \textbf{\texttt{MainPage.js}}, and the scattered local storage reads were all accepted this way, repeatedly, before any of them were treated as a problem. Oversight only became substantive once an issue was obvious enough to notice on its own --- the file had already grown past 180 lines, or local storage was already being read independently in three separate places --- at which point the fix required undoing an accumulation rather than preventing one decision. Had the same level of scrutiny been applied earlier and consistently, at the point each individual decision was made rather than only once its consequences had piled up, these problems would most plausibly have been caught and corrected while still cheap to fix.

The third is that generating code which is technically valid is not the same as generating code that makes good decisions for its context. The clearest instance of this is the password-hashing performance problem: the initial implementation used Werkzeug's default hashing function, which is deliberately expensive to compute as a security measure, and this default is a reasonable choice in general. But applied without regard to the local development server's single-threaded, synchronous request handling, it introduced a delay noticeable enough to be flagged during testing. The fix --- switching to bcrypt, where the hashing cost can be tuned explicitly --- did not require abandoning the underlying security principle, only adjusting it to fit the actual environment the code was running in. The same pattern, in a different form, runs through the architectural problems already discussed: a service layer built from scratch for each component, a layout component that kept absorbing new responsibilities, and authentication state read independently by whichever component needed it were all technically functional at the moment each was written, and none of them represented, in isolation, a wrong decision. What they lacked was any consideration of whether the choice would still hold up as the application grew around it --- which is a different kind of failure than producing code that simply does not work.

